\documentclass[10pt,a4paper]{article}
\usepackage[left=2cm, right=2cm]{geometry}
\setlength{\parindent}{0pt}
\usepackage{parskip}
\usepackage[super]{nth}

\begin{document}


{\centering \Large Computer Science Tripos -- Part II -- Progress Report\par}
\vspace{1cm}
{\centering \Huge Next Generation CHERI Tag Controller\par}
\vspace{1.5cm}
{\centering \large Kofi Wilkinson, Homerton College\par}
{\centering \large ksw40@cam.ac.uk\par}
{\centering \large Project Originator: Dr Jonathan Woodruff\par}
{\centering \large \nth{2} February 2024\par}
\vspace{1.5cm}

%\begin{tabular}{ll}
%\textbf{Project Supervisor:} & Dr Jonathan Woodruff\\
%\textbf{Director of Studies:} & Dr John Fawcett\\
%\textbf{Project Checkers:} & Prof. Alan Blackwell \& Prof. Srinivasan Keshav
%\end{tabular}

\textbf{Project Supervisor:} Dr Jonathan Woodruff\par
\textbf{Director of Studies:} Dr John Fawcett\par
\textbf{Project Checkers:} Prof. Alan Blackwell \& Prof. Srinivasan Keshav
\vspace{2cm}


\section*{Description of Progress}
So far, I have obtained a trace file and successfully written a decoder in C++. I have also written C++ that parses the contents of the
file (which contains the DRAM access trace of a particular workload), and simulates
the execution of a given tag controller algorithm running on said workload. This
produces an output that can be consumed by ChampSim, an existing C++ cache
simulator, which currently provides me with some basic metrics, such as tag cache hit and miss frequencies.
I have implemented the hierarchical tag controller algorithm detailed in the
paper \textit{Efficient Tagged Memory}, and have run successful simulations with
it. This code however is currently untested, and there are many opportunities for errors
to have slipped in.

\section*{Revisiting the Schedule}
\subsection*{Status}
I am more than halfway through the core of my project -- these tasks remain:
\begin{itemize}
	\item Gain access to the traces on the CL network
	\item Select a useful set of analysis/visualisation tools
	for the simulator, and verify with my supervisor
	\item Implement these tools
	\item Implement Arm's Morello controller algorithm
	\item Test both algorithms and evaluate them using the tools, recording
		my findings in a formal document.
\end{itemize}

I have not yet begun work on any of the exetnsions of my project.

\subsection*{Difficulties}
One great unforseen difficulty was how long decoding successfully would take.
A lack of
documentation required me to
look through some code of an alumnus who once worked with
these files, to engage in a long email exchange with said alumnus, and demanded
much trial and error.

Another difficulty was finding the correct output trace format
my program should use so that ChampSim can parse them.
ChampSim has no documentation
regarding its trace formats. Again, this had me looking through
(ChampSim's) code.

A further difficulty was that my Michaelmas was relatively heavy
with lecture courses. I initially took 5 and dropped one
after spending some time on it. I did not originally plan this, and my
intention only changed after learning more about this year's courses. Thus,
my original workplan did not factor this in, leading to me being further
behind than expected.

\subsection*{Adapted Workplan}
I am taking 2 courses in Lent and 1 in Easter. This leaves me with more time to catch up.

Gaining access to the traces should be straightforward via communications with
CL staff.
I will implement the Morello algorithm next. After, I will brainstorm some
ideas for metrics that the analysis/visualisation tools could provide, and will
implement them. This should be done by 07/02.\\
I will then test the algorithms, evaluate them with the tools, and record my
findings. This should be done by 12/02.\\
This leaves me 4 weeks behind my original schedule. With a total of 3 weeks
of slack period remaining in the original schedule (not counting my 2 week
safety-net slack period at the end), and this term's lighter lecture load, I should
be able to make up for the remaining 1 week task debt.

\end{document}
